\documentclass[12px]{article}

\title{Lezione 4 Metodi Numerici}
\date{2026-03-10}
\author{Federico De Sisti}

\input{../../setup.tex}

\begin{document}
\maketitle
	\newpage
	\subsection{Primi schemi}
	usando $f_n = f(t_n,u_n)$\\
	$FE\ \ \ u_{n+1} = u_n + hf_n$\\
	$BE \ \ \ u_{n+1} = u_n + hf_{n+1}$\\
	$CN \ \ \ u_{n+1} = u_n + \frac h 2 (f_n  + f_{n+1})$\\
	$H \ \ \ \ \ u_{n+1} = u_n + \frac h 2 (f_n + f(t_{n+1}, u_n + h f_n))$\\
	Tutti questi metodi hanno una filosofia comune
	\subsection{Filosofia}
	\[
		(ODE) \begin{cases}
		\dot y(t) = f(t,y(t)) \ t\in (0,T)\\
		y(0) = u_0
	\end{cases} \Leftrightarrow (V)\ y(t) = y_0 + \int_0^t f(s,y(s))ds
	.\] 
	Chiaramente a patto di avere una soluzione abbastanza regolare, se la dinamica è continua e anche la mia soluzione lo è, basterebbe.\\
	Chiamiamo la forma dell'ODE forma forte, e la forma di Volterra forma debole.\\

	Nel caso dell'ODE, cosa posso approssimare? Dato che è tutto dato posso solamente approssimare la derivata $\dot y$. (differenze finite)\\

	Nel caso di Volterra, posso approssimare invece l'integrale. (formule di quadratura)\\

	Ricordando che $\dot y = \lim_{\delta \rightarrow 0} \frac{y(t+\delta) - y(t)}{\delta} \\
	\Rightarrow  \phi(\delta) \xrightarrow{\delta \rightarrow 0} 0 \Rightarrow  \dot y(t) = \lim_{\delta \rightarrow 0} \frac{y(t + \phi(\delta)) - y(t) }{\phi(\delta)}$ è una formulazione equivalente\\
	
	Proviamo a fare una scelta particolare $\delta = h$, ovvero il passo di discretizzazione (quanti nodi ci sono)\\
	 \[
		 \dot y(t) = \lim_{h \rightarrow 0}\frac{y(t+h)-y(t)} h\approx \frac{y(t+h) - y(t)}h
	.\] 
	\textbf{Esempio}\\
	$\lim_{x \rightarrow 0}\frac{\sin x}x$ in matlab se sostituissi la $x$, verrebbe Nan, per ovviare a questo problema possiamo fare lo sviluppo di Taylor
	\[
		\Rightarrow  f(x) \sim f(x_0) + f'(x_0)(x-x_0) + \frac{f''(x_0)} 2 (x-x_0)^2 + \ldots
	.\] 
	$ \Rightarrow  \sin x \sim 0 + x - \frac{x^3}6 + \ldots$ \\
	\[
		\frac{\sin x}x \sim \frac{x - \frac {x^3}6 + O(x^3)}x = 1 - \frac{x^2}6 + o(x^2) \xrightarrow{x \rightarrow 0} 1
	.\] 
	Quindi tornando al nostro rapporto incrementale con $h$ (necessariamente) fissato\\
	\[
		\frac{ y(t+h) - y(t)}h =\frac{ \cancel{y(t)} + y'(t) + \frac 12 h^2y''(\xi) - \cancel{y(t)}}{h}
	.\] 
	con $t < \xi < t+ h$\\
	\[\frac{y(t+h) - y(t)}h = y'(t) + \frac 12 h y''(\xi)\]
	Quindi se approssimo ho un errore proporzionale a $h$
	\[
		err = \left|\frac{y(t+h)-y(t)}h - y'(t) \right| = \frac 12h|y''(\xi)|\leq \frac{M h}2
	.\] 
	con $M = \max_{0 < \xi < h}|y''(\xi)|$ \\
	Se voglio che $\frac{Mh}2 < tol$  dovrò avere  $h < \frac {2 tol} M$\\
	Tornando alla nostra ODE
	\[
	\begin{cases}
		\dot y(t) = f(t,y(t))\\
		y(0) = y_0
	\end{cases} \rightarrow \frac{y(t+h) + y(t)}h = f(t,y(t))
	.\] 
	Leggiamo questa espressione con $t = t_n\ \ n = 0,\ldots, N$
	 \[
		 \frac{y(t_n + h) - y(t_n)}h = f(t_n,y(t_n))
	.\] 
	ovvero $y(t_{n+1}) = y(t_n) + hf(t_n,y(t_n))$, il metodo di Eulero in avanti\\

	Posso ora provare ad approssimare la derivata prima in un altro modo.
	 \[
		 \dot y(t) = \lim_{ h \rightarrow 0}\frac{ y(t) - y(t-h)}h
	.\] 
	\[
		\dot y(t) = f(t,y(t)) \rightarrow \frac{y(t) - y(t-h)}h = f(t,y(t))
	.\] 
	$y(t) = y(t-h) + h f(t,y(t))$\\
	Leggo in $t = t_{n+1} \Rightarrow  y(t_{n+1} = y(t) + hf(t_{n+1},y(t_{n+1}))$ ovvero Eulero all'indietro.\\
	Cambiamo punto di vista:\\
	$y(t) = y_0 + \int_0^tf(s,y(s))ds$\\
	Posso approssimare con dei quadrati $\int_0^t f(s,y(s))ds\sim tf(0,y(0))$\\
	Scriviamo ora la forma di volterra in $t+h$
	 \[
		 y(t+h) =y_0 + \int_0^{t+h}f(s,y(s))ds = y_0 + \int_0^tf(s,y(s))ds + \int_{t}^{t+h}f(s,y(s))ds
	.\] 
	se $y = t_n$
	 \[
		 y(t_n + h) = y(t_{n+1}) = y_0 + \int_0^{t_n}f(s,y(s))ds + \int_{t_n}^{t_{n+1}}f(s,y(s))ds
	\] 
	\[
	 = y(t_n) + \int_{t_n}^{t_{n+1}}f(s,y(s))ds
	.\] 
	\[
		y(t_{n+1}) = y(t_n) + \int_{t_n}^{t_{n+1}}f(s,y(s))ds
	.\] 
	Approssimazione con il metodo dei rettangoli (a sinistra)
	\[
		\int_{t_n}^{t_{n+1}} f(s,y(s))ds\sim h f(t_n,y(t_n))
	.\] 
	Ovvero Eulero in avanti\\
	Approssimazione con il metodo dei rettangoli (a destra)
	\[
		\int_{t_n}^{t_{n+1}}f(s,y(s))ds \sim hf(t_{n+1},y(t_{n+1}))
	.\] 
	Ovvero Eulero all'indietro\\
	Approssimazione con il metodo dei trapezi
	\[
		\int_{t_n}^{t_{n+1}}f(s,y(s))ds\sim \frac h 2 (f(t_n,y(t_n)) + f(t_{n+1},y(t_{n+1})))
	.\] 
	Ovvero Crank Nicholson\\
	Heun la vediamo dopo, non utilizza valori della griglia.\\
	Introduciamo il $\theta$-Metodo,\\
	Scelgo $\theta\in [0,1]$
	 \[
		 \int_{t_n}^{t_{n+1}}f(s,y(s))ds \sim h \left((1-\theta)f(t_n,y(t_n)) + \theta f(t_{n+1},y(t_{n+1})) \right)
	.\] 
	Per $\theta = \frac 12 $ ottengo Crank Nicholson.\\
	Cosa succede se
	 \[
	 	 \frac{y(t+h) - y(t-h)}{2h} = \frac{y(t + h) - y(t) + y(t) - y(t+h)}{2h} 
	\] 
	\[
		= \frac 12 \left(\frac{y(t+h)-y(t)}h \right) + \frac 12  \left( \frac{y(t) - y(t-h)}h \right) \xrightarrow{h \rightarrow 0} \frac 12\dot y(t) + \frac 12 \dot y(t) = \dot y(t)
	.\]
	Questo metodo ci permette di semplificare i termini di ordine superiore nello sviluppo di Taylor.\\
	Faccio lo sviluppo di Taylor di $y(t+h)$ in  $t$ \\
	\[
	y(t+h) = y(t) + \dot y(t)h + \frac 12 \ddot y(t)h^2 + \frac 16 \dddot y(\xi)h^3\ \ t < \xi < t + h
	.\] 
	\[
	y(t-h) = y(t) - \dot y(t)h + \frac 12 \ddot y(t)h^2  - \frac 16\dddot y(\eta)h^3\ \ \ t-h < \eta < t 
	.\] 
	Facendo la differenza rimuovo i termini pari
	\[
	y(t+h) - y(t-h) = \dot y(t)2h + \frac 16 h^3(\dddot y(\xi) + \dddot y(\eta))
	.\] 
	Dividendo entrambe le parti per $2h$
	 \[
		 \frac{y(t+h)-y(t-h)}{2h} = \dot y(t) + \frac{h^2}{12}(\dddot y(\xi) + \dddot y(\eta))
	.\] 
	Possiamo quindi essere più precisi, ma al costo di utilizzare dei nodi diversi, utilizza non solo $t_{n}$ ma anche  $t_{n-1}$, ovvero è uno schema a 2 passi.\\
	$\dot y(t) = f(t,y(t)) \Rightarrow  \frac{y(t+h) - y(t-h)}{2h} = f(t,y(t))$ \\
	$t=t_n$
	 \[
		 \frac{y(t_{n+1}) - y(t_{n-1})}{2h} = f(t_{n},y(t_{n})) \Rightarrow  y(t_{n+1}) = y(t_{n-1}) + 2h f(t_n,y(t_n))
	.\] 
	Questo è un metodo esplicito a 2 passi, viene chiamato il mid-point, poichè dal punto di vista integrale è come se approssimassimo gli integrali con li quadrato di base il punto medio\\
	\[
		y(t+h) = y(t) + \int_{t_n}^{t_{n+1}}f(s,y(s))ds = y(t-h) + \int_{t-h}^tf(s,y(s))ds + \int_t^{t+h}f(s,y(s))ds
	\] \[
	= y(t-h) + \int_{t-h}^{t+h}f(s,y(s))ds
	.\] 
	\subsection{Analisi dei metodi ad un passo (espliciti)}
\[
	u_{n+1} = u_n + h\Phi(t_n,u_n; h, f)
.\] 
il punto e virgola è per dire che in generale deve dipendere da $h$ e dalla dinamica del problema.\\
$\Phi$ è la funzione incremento.\\
 \textbf{Esempi}\\
 $\Phi(t_n,u_n) = f(t_n,u_n)$ \hfill  (FE)\\
 $\Phi(t_n,u_n) = \frac 12 (f(t_n,u_n) + f(t_{n+1},u_n + hf(t_n,u_n)))$ \hfill (H)
 \subsubsection{Consistenza}
Se sostituisco la soluzione esatta di ODE nello schema commetto un errore
\[
	y(t_{n+1}) = y(t_n) + h\Phi(t_n,y(t_n))+ \varepsilon_{n+1}
.\] 
o come formulazione alternativa
\[
	\varepsilon_{n+1}:= y(t_{n+1})-y(t_n) -h \Phi(t_n,y(t_n))
.\] 
Chiamo $\varepsilon_{n+1}$ residuo al tempo  $t_{n+1}$\\
Riscrivo $\varepsilon_{n+1} = hz_{n+1}(h)$ LTE (Local Truncation Error)

\end{document}
